\documentclass[10pt,a4paper]{article}
\usepackage[margin=0.7in]{geometry}
\usepackage{amsmath,booktabs,listings}
\usepackage[hidelinks]{hyperref}
\setlength{\parindent}{0pt}
\setlength{\parskip}{4pt}
\lstset{basicstyle=\footnotesize\ttfamily,breaklines=true,columns=fullflexible,
        keepspaces=true,showstringspaces=false,aboveskip=5pt,belowskip=5pt}
\newcommand{\task}[1]{\medskip\textbf{#1}\par}
\begin{document}
\begin{center}
{\Large\bfseries Lab 4: PySpark}\\[3pt]
26AI5703 -- Data Engineering \& Big Data Systems Lab
\end{center}
\textbf{Aim:} Process the supplied datasets using Spark RDDs.
The code and results are presented below in task order.

\task{1a. Spark installation and test}
Spark 4.2.0 and Java 21 were installed. The Spark shell was started with:
\begin{lstlisting}
spark-shell --master 'local[2]'
\end{lstlisting}
Shell command: \texttt{sc.parallelize(Seq(1, 2, 3)).sum()}. Output: \texttt{6.0}.
The equivalent Python check and setup are:
\begin{lstlisting}
import findspark
findspark.init()

from operator import add
from pyspark import SparkConf, SparkContext

conf = (SparkConf().setMaster("local[2]").setAppName("Lab4")
        .set("spark.hadoop.fs.defaultFS", "file:///"))  # Read local files, not HDFS.
sc = SparkContext.getOrCreate(conf)
sc.setLogLevel("ERROR")
print("Spark:", sc.version, "| Test sum:", sc.parallelize([1, 2, 3]).sum())
\end{lstlisting}
\textbf{Output:} \begin{lstlisting}
Spark: 4.2.0 | Test sum: 6
\end{lstlisting}

\task{1b. Word count}
\texttt{flatMap} splits lines into words. Each word becomes \texttt{(word, 1)},
and \texttt{reduceByKey(add)} adds the counts of repeated words.
The sample file contains \texttt{hello spark}, \texttt{hello hadoop}, and \texttt{spark spark}
on three lines. The earlier Hadoop input was not supplied.
\begin{lstlisting}
word_counts = (sc.textFile("data/wordcount.txt").flatMap(str.split)
               .map(lambda word: (word, 1)).reduceByKey(add)
               .sortByKey().collect())
print(word_counts)
\end{lstlisting}
\textbf{Output:}
\begin{center}
\begin{tabular}{lr}
\toprule
Word & Count \\
\midrule
hadoop & 1 \\
hello & 2 \\
spark & 3 \\
\bottomrule
\end{tabular}
\end{center}

\task{1c. Rating distribution for every movie}
MovieLens 100K has 100,000 ratings for 1,682 movies.
Count each \texttt{(movieID, rating)} pair, then collect the five counts per movie.
Missing ratings have count zero.
\begin{lstlisting}
ratings = sc.textFile("data/ml-100k/u.data").map(str.split)
rating_counts = (ratings.map(lambda r: ((int(r[1]), int(r[2])), 1))
                 .reduceByKey(add))
movie_counts = (rating_counts.map(lambda x: (x[0][0], (x[0][1], x[1])))
                .groupByKey().mapValues(dict).sortByKey().collect())
movie_distribution = [(movie, [counts.get(rating, 0) for rating in range(1, 6)])
                      for movie, counts in movie_counts]
print(movie_distribution[:10])
\end{lstlisting}
\textbf{Output, first five movies (from the ten-row preview):}
\begin{center}
\begin{tabular}{rrrrrr}
\toprule
Movie ID & 1 star & 2 stars & 3 stars & 4 stars & 5 stars \\
\midrule
1 & 8 & 27 & 96 & 202 & 119 \\
2 & 8 & 17 & 55 & 42 & 9 \\
3 & 11 & 20 & 25 & 23 & 11 \\
4 & 6 & 24 & 57 & 93 & 29 \\
5 & 4 & 11 & 32 & 33 & 6 \\
\bottomrule
\end{tabular}
\end{center}
\newpage
\task{2. Average number of friends for every age}
Input columns: ID, name, age, and number of friends (no header).
The dataset contains 500 people across 52 ages.
\texttt{mapValues} changes only the value of a key--value pair; the age key stays the same.
It first changes a friend count to \texttt{(count, 1)}. \texttt{reduceByKey} adds
both parts for each age, and the second \texttt{mapValues} calculates:
\[
\text{average friends at an age} =
\frac{\text{total friends of people of that age}}{\text{number of people of that age}}.
\]
\begin{lstlisting}
friends = sc.textFile("data/friends_test.csv").map(lambda line: line.split(","))
average_friends = (friends.map(lambda r: (int(r[2]), int(r[3])))
                   .mapValues(lambda count: (count, 1))
                   .reduceByKey(lambda a, b: (a[0] + b[0], a[1] + b[1]))
                   .mapValues(lambda total: total[0] / total[1])
                   .sortByKey().collect())
for age, average in average_friends:
    print(age, round(average, 2))
\end{lstlisting}
\textbf{Output: all ages, rounded to two decimal places.}
\begin{center}
\begin{tabular}{rr@{\qquad}rr@{\qquad}rr}
\toprule
Age & Average & Age & Average & Age & Average \\
\midrule
18 & 343.38 & 36 & 246.60 & 54 & 278.08 \\
19 & 213.27 & 37 & 249.33 & 55 & 295.54 \\
20 & 165.00 & 38 & 193.53 & 56 & 306.67 \\
21 & 350.88 & 39 & 169.29 & 57 & 258.83 \\
22 & 206.43 & 40 & 250.82 & 58 & 116.55 \\
23 & 246.30 & 41 & 268.56 & 59 & 220.00 \\
24 & 233.80 & 42 & 303.50 & 60 & 202.71 \\
25 & 197.45 & 43 & 230.57 & 61 & 256.22 \\
26 & 242.06 & 44 & 282.17 & 62 & 220.77 \\
27 & 228.12 & 45 & 309.54 & 63 & 384.00 \\
28 & 209.10 & 46 & 223.69 & 64 & 281.33 \\
29 & 215.92 & 47 & 233.22 & 65 & 298.20 \\
30 & 235.82 & 48 & 281.40 & 66 & 276.44 \\
31 & 267.25 & 49 & 184.67 & 67 & 214.62 \\
32 & 207.91 & 50 & 254.60 & 68 & 269.60 \\
33 & 325.33 & 51 & 302.14 & 69 & 235.20 \\
34 & 245.50 & 52 & 340.64 &  &  \\
35 & 211.62 & 53 & 222.86 &  &  \\
\bottomrule
\end{tabular}
\end{center}
For example, people aged 18 have an average of 343.38 friends.
This is an average across those people, not the friend count of one person.

\newpage
\task{3 and 4. Minimum and maximum temperatures}
The CSV columns are \texttt{itemID, stationID, desc, temp}.
Skip the header and convert \texttt{temp} to an integer.
\texttt{filter} keeps only the requested measurement type:
\texttt{TMIN} for minimum temperatures, or \texttt{TMAX} for maximum temperatures.
Precipitation records are excluded.
\begin{lstlisting}
lines = sc.textFile("data/temp.csv")
header = lines.first()
temperatures = (lines.filter(lambda line: line != header)
                .map(lambda line: line.split(","))
                .map(lambda r: (r[0], r[1], r[2], int(r[3]))))

def temperature_extrema(kind, operation):
    selected = temperatures.filter(lambda r: r[2] == kind)
    overall = selected.map(lambda r: r[3]).reduce(operation)
    by_item = selected.map(lambda r: (r[0], r[3])).reduceByKey(operation).sortByKey().collect()
    by_station = selected.map(lambda r: (r[1], r[3])).reduceByKey(operation).sortByKey().collect()
    print("Overall:", overall)
    print("By item:", by_item)
    print("By station:", by_station)
    return overall, by_item, by_station

minimum_temperatures = temperature_extrema("TMIN", min)
maximum_temperatures = temperature_extrema("TMAX", max)
\end{lstlisting}
\texttt{reduce} combines all selected temperatures into one result.
\texttt{reduceByKey} combines only temperatures with the same item or station key.
The shared function uses \texttt{min} for task 3 and \texttt{max} for task 4.

\textbf{Output: overall and every item with temperature records.}
\begin{center}
\begin{tabular}{lrr}
\toprule
Group & Minimum (TMIN) & Maximum (TMAX) \\
\midrule
Overall & -148 & 323 \\
\texttt{EZE00100082} & -135 & 323 \\
\texttt{ITE00100554} & -148 & 323 \\
\bottomrule
\end{tabular}
\end{center}
\textbf{Output: first five station keys, out of 365.}
\begin{center}
\begin{tabular}{lrr}
\toprule
Station ID & Minimum (TMIN) & Maximum (TMAX) \\
\midrule
18000101 & -148 & -75 \\
18000102 & -130 & -44 \\
18000103 & -73 & -10 \\
18000104 & -74 & 0 \\
18000105 & -58 & 10 \\
\bottomrule
\end{tabular}
\end{center}
Temperatures retain the supplied numeric units. Grouping follows the CSV headers,
including the date-like values in \texttt{stationID}.

Stop Spark after completing all tasks:
\begin{lstlisting}
sc.stop()
\end{lstlisting}
\textbf{Result:} All tasks ran successfully. Outputs matched independent Python
calculations; all movie rating counts together total 100,000.

\textbf{Source:} \href{https://grouplens.org/datasets/movielens/100k/}{GroupLens MovieLens 100K}.
The accompanying code PDF contains the notebook cells and saved outputs.
\end{document}
